\section{Estruturas tectônicas e halocinese}

As estruturas tectônicas modificam a geometria das camadas e, por isso, influenciam a leitura do registro estratigráfico, a geração de relevo e a organização das bacias.

\subsection{Dobras}
\begin{definicao}

Os dobramentos \textbf{sinclinais} são estruturas geológicas em que as camadas rochosas se curvam de modo côncavo para cima, formando uma dobra cuja parte central contém, em geral, as camadas mais recentes. Em um sinclinal, os flancos mergulham em direção ao eixo da dobra.

\end{definicao}

\subsection{Falhas}

Uma \textbf{falha} é uma descontinuidade ao longo da qual ocorreu deslocamento relativo entre blocos rochosos. A classificação cinemática considera a direção do movimento:

\begin{itemize}
  \item \textbf{falha normal}: o bloco de teto desce em relação ao bloco de muro; associa-se principalmente à extensão;
  \item \textbf{falha reversa}: o bloco de teto sobe em relação ao bloco de muro; associa-se principalmente à compressão;
  \item \textbf{falha transcorrente}: o deslocamento predominante é horizontal e pode ser dextral ou sinistral;
  \item \textbf{falha oblíqua}: combina componentes vertical e horizontal importantes.
\end{itemize}

Falhas lístricas possuem superfície curva, geralmente com diminuição do mergulho em profundidade. Em bacias extensionais, falhas normais podem delimitar grábens e semigrábens. Quando o deslocamento ocorre simultaneamente à deposição e produz variação sistemática da espessura dos estratos, usa-se o termo \textbf{falha de crescimento} \parencite{fossen2016}.

\subsection{Exemplo: Bacia do Recôncavo}

A Bacia do Recôncavo integra o sistema rifte Recôncavo--Tucano--Jatobá, desenvolvido durante a ruptura continental relacionada à abertura do Atlântico Sul. Seu arcabouço apresenta blocos assimétricos controlados predominantemente por falhas normais. A Falha de Salvador limita a bacia a leste, enquanto a Falha de Maragogipe integra seu limite ocidental. O preenchimento registra sistemas continentais, especialmente lacustres, fluviais e deltaicos \parencite{anp_reconcavo,magnavita2005}.

Esse exemplo mostra a relação entre:

\[
\begin{gathered}
\text{extensão crustal}
\rightarrow
\text{falhamento normal}
\rightarrow
\text{subsidência}
\\[-0.2em]
\rightarrow
\text{criação de acomodação}
\rightarrow
\text{preenchimento sedimentar}.
\end{gathered}
\]

\subsection{Exemplo: estruturas do estado do Rio de Janeiro}

O mapa geológico estadual registra estruturas de diferentes naturezas e idades, incluindo falhas, zonas de cisalhamento e lineamentos que afetam o embasamento e as coberturas sedimentares. Sua interpretação exige distinguir estruturas antigas do Cinturão Ribeira de reativações posteriores relacionadas à evolução da margem continental e do relevo do Sudeste. Assim, uma direção cartográfica semelhante não prova, por si só, que duas falhas tenham a mesma idade ou cinemática \parencite{heilbron2016,fossen2016}.

\begin{definicao}
 Os dobramentos \textbf{anticlinais} são estruturas geológicas em que as camadas rochosas se curvam de modo convexo para cima, formando uma dobra cuja parte central contém, em geral, as camadas mais antigas. Em um anticlinal, os flancos mergulham em sentido oposto ao eixo da dobra.   
\end{definicao}

\subsection{Diápiros e estruturas salinas}

\begin{definicao}
Os \textbf{diápiros} são estruturas geológicas formadas pela ascensão de materiais mais dúcteis e menos densos através de camadas rochosas superiores mais rígidas e mais densas. Essa ascensão ocorre em razão da diferença de densidade entre os materiais e da capacidade de deformação plástica da massa ascendente.

Os \textbf{diápiros de sal} são diápiros constituídos predominantemente por evaporitos, especialmente halita, podendo também incluir anidrita, gipsita e outros sais. Em profundidade, o sal comporta-se de maneira dúctil e, por ser menos denso que grande parte das rochas sedimentares que o recobrem, tende a migrar para níveis superiores, deformando as camadas encaixantes.

Os diápiros de sal ocorrem em \textbf{bacias sedimentares evaporíticas}, isto é, em bacias nas quais houve deposição de espessas camadas de sal em ambientes de intensa evaporação. Sua formação inicia-se com a deposição desses evaporitos, seguida pelo soterramento por sedimentos mais densos. Com o aumento da sobrecarga, o sal passa a fluir lentamente devido ao seu comportamento plástico e à diferença de densidade, podendo ascender e formar estruturas como domos e muralhas salinas. Essas estruturas são comuns em regiões como o Golfo do México, o Mar do Norte, o Oriente Médio e a margem continental brasileira, especialmente em áreas associadas ao pré-sal \parencite{jacksonhudec2017}.

\end{definicao}
